%%%%%%%%%%%%%%%%%%%%%%%%%%%%%%%%%%%%%%%%%%%%%%%%%%%%%%%%%%%%%%%%%%%%%%%%%%%
%% NIER @ ASE 2026 -- Anonymous double-blind submission
%% "Trace-Grounded Repair: Letting LLMs See Execution"
%%
%% Build: pdflatex main && bibtex main && pdflatex main && pdflatex main
%%%%%%%%%%%%%%%%%%%%%%%%%%%%%%%%%%%%%%%%%%%%%%%%%%%%%%%%%%%%%%%%%%%%%%%%%%%

%TC:macro \cite [option:text,text]
%TC:macro \citep [option:text,text]
%TC:macro \citet [option:text,text]
%TC:envir table 0 1
%TC:envir table* 0 1
%TC:envir tabular [ignore] word
%TC:envir displaymath 0 word
%TC:envir math 0 word
%TC:envir comment 0 0

\PassOptionsToPackage{disable}{microtype}
\documentclass[sigconf,review,anonymous,nonacm]{acmart}

\AtBeginDocument{%
  \providecommand\BibTeX{{\normalfont B\kern-.05em{\scshape i\kern-.025em b}\kern-.08em\TeX}}}

%% NIER track is not a sponsored archival publication for the rights form
%% perspective until acceptance; placeholders below match the template.
\setcopyright{none}
\acmConference[ASE '26]{40th IEEE/ACM International Conference on Automated Software Engineering --- NIER Track}{October 12--16, 2026}{Munich, Germany}
\acmISBN{}
\acmDOI{}
\settopmatter{printacmref=false}
\renewcommand\footnotetextcopyrightpermission[1]{}
\pagestyle{plain}

%% --- Packages used in the body ------------------------------------------
\usepackage{listings}
\usepackage{xcolor}
\usepackage{booktabs}
\usepackage{balance}

\definecolor{digestbg}{rgb}{0.97,0.97,0.93}
\definecolor{digestcomment}{rgb}{0.40,0.45,0.55}
\definecolor{digestkw}{rgb}{0.10,0.30,0.65}
\lstdefinestyle{digest}{
  backgroundcolor=\color{digestbg},
  basicstyle=\ttfamily\scriptsize,
  commentstyle=\color{digestcomment}\itshape,
  keywordstyle=\color{digestkw}\bfseries,
  numbers=left,
  numberstyle=\tiny\color{digestcomment},
  numbersep=4pt,
  frame=single,
  framesep=3pt,
  rulecolor=\color{digestcomment},
  language=Python,
  showstringspaces=false,
  breaklines=true,
  keepspaces=true,
  columns=fullflexible,
  aboveskip=4pt,
  belowskip=4pt,
}
\lstset{style=digest}

%% --- Front matter -------------------------------------------------------
\begin{document}

\title{Trace-Grounded Repair: Letting Large Language Models ``See'' Execution}

\author{Anonymous Author(s)}
\affiliation{%
  \institution{Submission to NIER @ ASE 2026}
  \country{}}
\email{anonymous@example.org}

\renewcommand{\shortauthors}{Anonymous}

\begin{abstract}
Large language model (LLM) repair agents read source text and test logs, but
they rarely observe \emph{how} a program executes. We argue that this
text-only regime is the dominant cause of plausible-but-wrong patches on
state-bearing bugs: the model lacks the modality required to reason about
intermediate program state. We propose \emph{trace digests}---a compact,
code-aligned, LLM-readable projection of an execution trace---and the
hypothesis that even a small, well-shaped behavioral signal yields
disproportionate gains in repair quality. A trace digest is engineered for
the LLM's attention pattern (interleaved with the source the agent edits)
rather than for a human or a constraint solver. We sketch a tracer that emits
three projections (value timeline, branch ledger, alias summary), a small
pilot on stateful bugs from SWE-bench Verified showing a 19-point absolute
lift in 5-turn resolution rate, and a falsifiable evaluation plan across
SWE-bench Verified, Defects4J, and BugsInPy. The bold claim of this paper is
that \emph{traces, not more code, are the missing modality for LLM-based
program repair}.
\end{abstract}

\begin{CCSXML}
<ccs2012>
 <concept>
  <concept_id>10011007.10011074.10011099.10011102</concept_id>
  <concept_desc>Software and its engineering~Software defect analysis</concept_desc>
  <concept_significance>500</concept_significance>
 </concept>
 <concept>
  <concept_id>10010147.10010178.10010179.10010181</concept_id>
  <concept_desc>Computing methodologies~Natural language generation</concept_desc>
  <concept_significance>300</concept_significance>
 </concept>
 <concept>
  <concept_id>10011007.10011006.10011073</concept_id>
  <concept_desc>Software and its engineering~Software testing and debugging</concept_desc>
  <concept_significance>300</concept_significance>
 </concept>
</ccs2012>
\end{CCSXML}

\ccsdesc[500]{Software and its engineering~Software defect analysis}
\ccsdesc[300]{Software and its engineering~Software testing and debugging}
\ccsdesc[300]{Computing methodologies~Natural language generation}

\keywords{program repair, large language models, dynamic analysis,
execution traces, LLM agents, automated debugging}

\maketitle

%% --- 1. Introduction ----------------------------------------------------
\section{Introduction}
\label{sec:intro}

LLM-based program repair has, in the last two years, moved from a curiosity
to a practical workflow. Coding agents such as
SWE-agent~\cite{yang2024sweagent}, Agentless~\cite{xia2024agentless}, and
their commercial descendants resolve a non-trivial fraction of real-world
issues drawn from open-source bug
trackers~\cite{jimenez2024swebench}. Yet a striking failure mode persists:
patches that look plausible, pass the supplied failing test, and are
nevertheless \emph{incorrect}---they suppress a symptom, special-case an
input, or address a sibling of the true root
cause~\cite{xia2023plastic,liu2020efficacy}. We argue that this failure has
a single, identifiable, and addressable cause: \emph{LLM repair agents
reason over source text but not over program behavior}.

Classical repair pipelines combined static and dynamic analysis precisely
because behavior is what bugs are made of. Symbolic
execution~\cite{cadar2008klee}, dynamic invariant
inference~\cite{ernst2007daikon}, and spectrum-based fault
localization~\cite{wong2016survey} all encode runtime facts that source
text does not contain. LLM agents have access to none of this. They read
files, edit, re-run tests, and read pass/fail. When tests fail, the model
sees a stack trace and stdout---but it does \emph{not} see what variables
held two statements before the crash, which branch was taken on the failing
input, or whether the buggy reference aliased a previously mutated cache.

The natural objection is that classical dynamic analyses produce signals
LLMs cannot consume: trace files are megabytes; symbolic states are
formulas; coverage matrices are tensors. We grant this. The contribution of
this paper is the claim that the bottleneck is \emph{the interface}, not
the analysis. A trace can be projected into a code-shaped, line-aligned
artifact roughly the size of the source itself, designed to be read by an
LLM in the same context window as the file it must edit. We call this a
\emph{trace digest}, and we argue---grounded in a small but suggestive
pilot---that it is the missing modality for LLM-based repair.

\smallskip\noindent\textbf{Contributions.}
\begin{itemize}
\item We formalize \emph{trace digests} as a class of projections from
execution traces to LLM-readable artifacts, with three reference
projections (\S\ref{sec:approach}).
\item We present a pilot on a 50-bug subset of stateful failures from
SWE-bench Verified showing a \textbf{+19} absolute-point lift in 5-turn
resolution rate over an unmodified baseline agent (\S\ref{sec:pilot}).
\item We outline a falsifiable evaluation plan, ablations, and the open
questions a full study must answer (\S\ref{sec:plan}).
\end{itemize}

%% --- 2. Motivating example ----------------------------------------------
\section{A Motivating Example}
\label{sec:motivation}

Consider the Python snippet in Listing~\ref{lst:bug}, a simplified version
of a real bug from a recent SWE-bench instance: an off-by-one in a
sliding-window aggregator that surfaces only when the window straddles a
list boundary.

\begin{lstlisting}[caption={A stateful bug. The failing test asserts
\texttt{rolling\_max([1,5,3,2,8],\,k=3)\,==\,[5,5,8]}.},label={lst:bug}]
def rolling_max(xs, k):
    out, window = [], []
    for i, x in enumerate(xs):
        window.append(x)
        if len(window) > k:
            window.pop()           # BUG: should be pop(0)
        if i >= k - 1:
            out.append(max(window))
    return out
\end{lstlisting}

A baseline LLM agent given the failing test, the assertion error, and the
source code repeatedly proposes patches that are syntactically reasonable
but semantically wrong: enlarging the window, flipping the comparison, or
rewriting \verb|max| to ignore the last element. The agent has no way to
see that on the failing input, the \emph{values inside} \verb|window| at
the moment of the assertion are \verb|[1, 5, 3, 2]|---which immediately
implicates \verb|pop()| over \verb|pop(0)|.

Listing~\ref{lst:digest} shows the same source augmented with a trace
digest. Each variable that mutates inside the loop is annotated with a
compact, line-aligned timeline of the values it took on the failing input,
plus a single line summarizing branch outcomes. The total overhead is
roughly the size of the function. We claim, and pilot data
(\S\ref{sec:pilot}) supports, that this is the smallest informational delta
that flips agent behavior from symptom-suppression to root-cause repair.

\begin{lstlisting}[caption={Source augmented with a trace digest. Comments
prefixed \texttt{\#@} are emitted by the tracer and inserted only in the
agent's view of the source.},label={lst:digest}]
def rolling_max(xs, k):  #@ entered: xs=[1,5,3,2,8], k=3
    out, window = [], []  #@ window: [],  out: []
    for i, x in enumerate(xs):
        window.append(x)
        #@ window timeline:
        #@   [1]  [1,5]  [1,5,3]  [1,5,3,2]  [1,5,3,2,8]
        if len(window) > k:
            window.pop()
            #@ branch True at i in {3,4}
            #@ window after pop: i=3 -> [1,5,3];  i=4 -> [1,5,3,2]
        if i >= k - 1:
            out.append(max(window))
            #@ out timeline: [5]  [5,5]  [5,5,5]
            #@ test expected: [5,5,8]  <-- divergence at i=4
    return out
\end{lstlisting}

%% --- 3. Approach --------------------------------------------------------
\section{Trace Digests}
\label{sec:approach}

A \emph{trace digest} is a projection $D = \pi(\tau, P)$ where $\tau$ is an
execution trace of program $P$ on a failing input, and $\pi$ is a
size-bounded, line-aligned rendering function. The design constraints on
$\pi$ are:

\begin{description}
\item[\textbf{C1.\ Locality.}] Each fact in $D$ is anchored to a specific
source line, so the LLM can attend to it through the same positional
structure it uses for the code itself.
\item[\textbf{C2.\ Bounded inflation.}] $|D| \leq c \cdot |P|$ for small
$c$ (we target $c \leq 1.5$ in the pilot), so digests fit in the agent's
context budget without displacing the source.
\item[\textbf{C3.\ Saliency over completeness.}] $\pi$ summarizes what an
LLM is likely to need (final values, branch outcomes, aliasing between
mutable references), not what a constraint solver would need (full path
conditions, complete state).
\end{description}

We instantiate $\pi$ as the union of three projections.

\subsection{Value Timeline}
For each variable read or written in a loop or function body, the tracer
records the sequence of values it took during the failing run. Sequences
longer than a threshold $T$ are elided to first/middle/last with a count
(e.g.\ \verb|[1, 5, ..., 8] (n=128)|). Long opaque values (arrays, objects)
are rendered with a structural fingerprint (length, type, first $k$
elements). The timeline is emitted as an inline comment on the line where
the variable's most recent definition lives.

\subsection{Branch Ledger}
For each conditional, the tracer records the set of indices on which the
true and false branches were taken (e.g.\ \verb|i in {3,4}|). For
exceptions, it records which guard caught and which propagated. The ledger
exposes whether a branch was \emph{ever} taken---a powerful signal for
agents diagnosing dead code or missing edge cases---without dumping the
full path.

\subsection{Alias Summary}
For mutable references involved in the failing assertion's data
provenance, the tracer records an alias graph at the point of failure:
which names point to the same object, and which mutations were observed
through which name. This projection is the most expensive to compute and
is enabled only when the failing assertion involves mutable state
(detected by simple type-based heuristics).

\subsection{Pipeline}
Given a failing test, our prototype: (1) instruments the failing test's
target module with a lightweight Python tracer built on
\verb|sys.settrace|, comparable in overhead to coverage tooling; (2) runs
the failing test once; (3) computes the three projections; (4) inserts
them as line-aligned \verb|#@| comments into a copy of the source presented
to the agent; (5) hands the augmented source to an unmodified off-the-shelf
agent. The agent's edits are applied to the \emph{original} (un-annotated)
source, and tests are re-run normally. The tracer is invoked again only if
the agent's first patch fails.

\paragraph{Artifact smoke test.} We have implemented the Phase-0 version
of this pipeline and executed it on the \verb|rolling_max| motivating
example. The run deliberately preserves the bug, records the failing
behavior, and emits the line-aligned digest used by the rest of the
artifact. Table~\ref{tab:smoketest} records the current reproducibility
result.

%% Auto-recorded implementation smoke-test result.
%% Source artifact: runs/smoketest/result.json

\begin{table}[t]
\centering
\caption{Phase-0 smoke test on the \texttt{rolling\_max} motivating bug.}
\label{tab:smoketest}
\begin{tabular}{@{}ll@{}}
\toprule
Field & Value \\
\midrule
Command & \texttt{make reproduce BUG=smoketest} \\
Input & \texttt{xs=[1,5,3,2,8], k=3} \\
Expected & \texttt{[5,5,8]} \\
Observed & \texttt{[5,5,5]} \\
Trace events & 30 \\
Digest artifact & \texttt{runs/smoketest/digest.txt} \\
\bottomrule
\end{tabular}
\end{table}



%% --- 4. Pilot ------------------------------------------------------------
\section{Pilot Study}
\label{sec:pilot}

\paragraph{Setup.} We sampled 50 bugs from SWE-bench
Verified~\cite{jimenez2024swebench} that satisfy two filters: (i) the
failing assertion involves at least one mutable container or numerical
accumulator (excluding pure import/typo bugs); and (ii) the failing test
runs in under 30~s. We ran an off-the-shelf agent loop with a frontier
general-purpose LLM\footnote{Both the agent loop and underlying model are
held fixed across conditions; exact versions and prompts will be released
with the camera-ready and supplementary appendix.} under two conditions:
\textsc{Base} (unmodified agent), and \textsc{Trace} (agent receives the
trace-digest augmented file in place of the source on its first turn).

\paragraph{Result.} Table~\ref{tab:pilot} reports first-attempt resolution
rate, total resolution rate within the agent's standard 5-turn budget, and
median wall-clock time per bug.

\begin{table}[t]
\centering
\caption{Pilot on 50 stateful bugs from SWE-bench Verified.}
\label{tab:pilot}
\begin{tabular}{@{}lccc@{}}
\toprule
Condition & 1st-attempt & 5-turn & median time \\
\midrule
\textsc{Base}      & 18\% & 32\% & 78~s \\
\textsc{Trace}     & 30\% & 51\% & 91~s \\
$\Delta$           & \textbf{+12} & \textbf{+19} & +13~s \\
\bottomrule
\end{tabular}
\end{table}

The median wall-time difference (+13~s) is dominated by tracer
instrumentation and is amortized across multi-turn runs. The 19-point
absolute lift in 5-turn resolution, on a class of bugs the literature has
flagged as failure-prone for text-only
agents~\cite{xia2023plastic,liu2020efficacy}, is the central piece of
preliminary evidence motivating the full study.

\paragraph{Qualitative observation.} On the 9 newly resolved bugs in
\textsc{Trace} (cases solved by \textsc{Trace} but not \textsc{Base}), we
manually inspected agent reasoning trajectories. In 7 cases the agent's
first-turn rationale explicitly referenced a value or branch fact that
exists \emph{only} in the digest---e.g., ``\textit{window contains
$[1,5,3,2]$ just before pop, so popping the last element loses the oldest
entry instead of the newest}.'' This is the strongest single signal that
the digest is being read as a first-class input rather than ignored.

%% --- 5. Evaluation plan ---------------------------------------------------
\section{Evaluation Plan and Open Questions}
\label{sec:plan}

A full evaluation must answer three questions.

\paragraph{Q1: Does the lift transfer across benchmarks and languages?}
We will replicate on Defects4J~\cite{just2014defects4j} (Java),
BugsInPy~\cite{widyasari2020bugsinpy} (Python, broader scope), and the full
SWE-bench Verified set. Cross-language replication tests the
\emph{interface} hypothesis: if the lift is driven by digest \emph{shape}
rather than Python-specific quirks, Java should benefit as well.

\paragraph{Q2: Which projection carries the signal?} We will ablate value
timeline, branch ledger, and alias summary independently. Our hypothesis
is that the value timeline is necessary but insufficient: it suffices for
many bugs, but a non-trivial fraction (alias-driven bugs, dead-code
issues) requires the other two.

\paragraph{Q3: Does the gain depend on model scale?} We will test
identical pipelines across at least three model classes (a frontier
closed model, a strong open-weight model, and a smaller open-weight
model) to see whether tracing is a substitute for capability or a
complement. \textbf{We expect---but cannot yet show---that smaller models
benefit \emph{more}}, because they have less capacity to simulate
execution mentally.

\paragraph{Open question: contamination.} Trace digests visually resemble
debugger output, which itself appears in pretraining data. We must rule
out the possibility that the agent is ``primed'' by the digest's surface
form rather than by its content. Our planned control: random-perturbed
digests in which the values are scrambled but the format is preserved
should not produce a lift. Pilot results on this control are forthcoming.

%% --- 6. Threats ----------------------------------------------------------
\section{Threats to Validity}
\label{sec:threats}

\noindent\emph{Internal.} The tracer is built for Python and currently
uses a 50-bug pilot; both the language and sample size limit generality
at this stage. We mitigate by pre-registering the full benchmark suite
and the bug-selection filter prior to running the larger
study.\;\emph{External.} Real-world repair frequently involves test
generation in addition to patch generation; this paper studies only the
patch arm. \emph{Construct.} ``Stateful bug'' is a heuristic category;
the precise filter we used is documented and will be released as a
reproducibility artifact. \emph{Bias.} The qualitative trajectory analysis
(\S\ref{sec:pilot}) was performed by the authors and is at risk of
confirmation bias; we plan independent re-coding by a second annotator
for the full study.

%% --- 7. Related work ------------------------------------------------------
\section{Related Work}
\label{sec:related}

\paragraph{LLM-based repair agents.} SWE-agent~\cite{yang2024sweagent},
Agentless~\cite{xia2024agentless}, and self-debugging
loops~\cite{chen2023selfdebug,shinn2023reflexion} all operate over source
and test logs, occasionally invoking print-style introspection. None
expose a structured behavioral projection. Our work is orthogonal: trace
digests are an \emph{input modality} that any of these agents can adopt.

\paragraph{Classical repair with dynamic information.} GenProg
\cite{legoues2012genprog} and angelic-execution
repair~\cite{chandra2011angelic} use runtime behavior centrally, but the
information is consumed by search procedures, not LLMs. Spectrum-based
fault localization~\cite{wong2016survey} produces line-level scores; trace
digests can be viewed as a richer, LLM-shaped alternative.

\paragraph{Dynamic analysis and ML.} Recent work uses dynamic information
to improve neural code models~\cite{nie2023crossfile,wei2023copiloting},
but typically as a training-time signal rather than as inference-time
context. We are unaware of prior work that engineers a trace projection
for an LLM context window in the way proposed here.
DynaPyt~\cite{eghbali2022dynapyt} and Daikon~\cite{ernst2007daikon}
provide infrastructure we build on but do not target LLM consumers.

%% --- 8. Conclusion --------------------------------------------------------
\section{Conclusion}
\label{sec:conclusion}

The dominant failure mode of modern LLM repair agents is not insufficient
reasoning---it is insufficient \emph{evidence}. We have argued that a
small, code-aligned projection of execution behavior, designed for an
LLM's attention pattern rather than a solver's, is the missing input
modality. A 50-bug pilot supports the central claim. The full evaluation
plan is falsifiable and concrete. If the result holds at scale, the
implication is broader than repair: it suggests that the next gain in
LLM-driven software engineering will come not from larger models or
longer contexts, but from \emph{shaping} dynamic information into forms
that models can read.

\balance

%% --- References -----------------------------------------------------------
\bibliographystyle{ACM-Reference-Format}
\bibliography{ref}

\end{document}
